\section{Discussion}
\label{sec:discussion}

\subsection{What the experiments establish}

The evidence identifies a useful algorithmic property: the same
factorized mixed-state construction produces competitive binary
assignments on sparse and dense MaxCut graphs and on two structurally
different planted Ising families. On the saved planted files, the
\QeFEM{} configuration has lower error than the declared \LQA{}
configuration in every instance, and it returns exact planted energies
in seven Wishart and 60 Chook trial populations. The G-set continuation
adds ten previously unattained reference cuts, eight of them recovered
with frozen settings on fresh seeds. These are concrete strengths, not
claims of quantum hardware advantage or universal optimality.

The dynamics help explain why those strengths are plausible. The
product-state expectation is cheap enough for a large replica batch;
high-temperature entropy gives a smooth starting regime; the driver
offers an additional local direction during part of the search; and
cooling lets the longitudinal coupling landscape determine the final
binary choices. Spectral scaling provides a principled coordinate for
the onset of longitudinal instability on zero-field graphs. None of
these mechanisms guarantees the best basin, but each gives a testable
design variable rather than an opaque postprocessing step.

The K2000 result imposes an equally clear boundary. A best cut of 33,293
is strong relative to the saved 33,249 \LQA{} cut, yet it remains 44
below the stored 33,337 reference. The best cut stops improving in the
recorded nested population sweep after 2,048 replicas. A larger
population alone therefore has no demonstrated route to reference
attainment in these data. The available \FEM{}/\dSB{} notebook outputs
reach 33,337 under their own protocols. The correct conclusion is that
the present \QeFEM{} K2000 configuration does not match those attained
cuts; the evidence does not establish that the state family is incapable
of matching them under a different schedule or update. Conversely,
theoretical inclusion of a classical endpoint does not prove a finite
\QeFEM{} run must outperform tuned \FEM{} dynamics. The endpoint
proposition in Sec.~\ref{sec:theory} makes that distinction exact.

\subsection{Why this can matter outside the present benchmarks}

Many scheduling, allocation, routing, and design tasks admit QUBO or
Ising formulations \citep{lucas2014}. Once such a formulation is
available, the variational state, entropy evaluation, readout, and
batch optimizer in \QeFEM{} do not depend on the application's name;
the objective enters through its fields and pair couplings. That
separation makes the method attractive for an optimization platform
that must support many binary models on accelerator hardware. The
experiments here support the \emph{technical possibility} of scaling a
single implementation from sparse G-set graphs to a 2,000-spin dense
graph and to planted suites with known targets. They do not yet measure
business value, deployment reliability, constraint satisfaction after
QUBO penalty tuning, or total energy and monetary cost on an industrial
workload. Those questions require application-specific experiments.

The distinction also shapes how performance should be communicated.
A result such as 60/110 planted-energy hits conveys something different
from an average percent error; together they show how often the method
fully solves a fixed file and how far the remaining runs miss. The
G-set count is compelling as an attainment inventory, but its adaptive
search budget must accompany it. On K2000, a 0.132\% reference-relative
residual sounds small, but the absolute 44-unit gap is the more direct
optimization fact on a signed-weight graph. Clear plots and complete
captions can present a strong algorithm without converting unlike
protocols into an unsupported ranking.

\subsection{Limitations and decisive next experiments}

The product ansatz omits intersite correlations and entanglement.
Low-temperature spin-glass landscapes can have many metastable basins,
and the optimizer follows a finite, schedule-dependent path rather than
the instantaneous variational minimum. The manual discrete-neighbour
update can further improve attainable cuts, but it is a surrogate
vector field, not the exact free-energy gradient. No general convergence
or approximation-ratio theorem is claimed for the complete practical
solver. In particular, an observed local spectral instability does not
imply global MaxCut optimality.

Several limitations are experimental. The plotted G-set total
accumulates adaptive tuning, historical runs, and fresh-seed
confirmation. K2000 has one graph and five canonical \QeFEM{} seeds;
the plotted planted series has one fixed file at each parameter.
Family-level hyperparameters were selected after exploration of related
data. The saved \LQA{} comparison evolves one state per trial, whereas
\QeFEM{} evolves 8,192, and its planted settings include a local driver
choice. No wall-clock-normalized or candidate-normalized ranking is
possible from these endpoints alone. Even the FEM/dSB K2000 notebook
is a contextual source rather than a matched three-way protocol; its
saved output predates a later device-selection code change, so the
exact numbers should be freshly reproduced before they become a formal
new head-to-head figure.

A decisive K2000 study would freeze graph input, precision, device,
scoring code, and a predeclared total-work budget across \QeFEM{},
\FEM{}, and \dSB{}. It would vary the \QeFEM{} schedule and the
G-set-derived manual-discrete update under separate, recorded
hyperparameter searches; then it would evaluate selected settings on
fresh seeds without moving the goalposts. Reporting per-replica cut
distributions, hit counts, time-to-target, and memory would determine
whether the remaining gap is a schedule issue, a basin-coverage issue,
or a more structural limitation. The current paper does not include
such a trial, and it does not extrapolate the G-set gains to K2000.

\section{Conclusion}

\QeFEM{} gives a mathematically explicit route from a binary Ising
objective to a differentiable, annealed search over local mixed states.
The quantum relative entropy supplies the free-energy principle; the
product ansatz makes it computationally tractable; the Bloch-field map
guarantees valid local states; and binary readout returns ordinary
feasible solutions. The theory locates a graph-dependent longitudinal
instability scale and shows exactly where the later discrete-neighbour
search departs from the smooth gradient.

Under the saved protocols, the method reaches 33 of 54 G-set reference
cuts across accumulated searches, improves the declared \LQA{}
comparison on every fixed Wishart and Chook instance, and gives a
33,293 K2000 cut while leaving a 44-unit reference gap. These results
make a coherent case for \QeFEM{} as a flexible classical optimization
framework with strong solution quality on several large binary
landscapes. They also identify the next falsifiable question: whether
the new search dynamics and carefully matched resources can close the
dense K2000 gap without relying on ever larger replica populations.

\section*{Author contributions}
R.D. conceived and developed the core \QeFEM{} method, implemented the
original solver, designed the analysis workflow, and drafted the
manuscript. F.P. contributed the nonlinear reparameterization,
optimizer and schedule refinement, and benchmark testing. Both authors
interpreted the results and take responsibility for the final text.
